\section{Introduction}
\label{sec:introduction}

Autonomous Large Language Model (LLM) agents are rapidly transitioning from conversational assistants into active operational entities capable of interacting with complex, external software environments. Frameworks such as LangGraph, AutoGen, CrewAI, and protocol standards like Anthropic's Model Context Protocol (MCP)~\cite{mcp2024spec} enable agents to dynamically invoke external APIs, query database systems, modify cloud infrastructure, and execute financial transactions. In these multi-tool orchestrations, an agent's control flow typically involves inspecting an external resource's state before triggering a state-modifying side effect based on the observation.

However, this decoupled check-and-act paradigm exposes a fundamental vulnerability: the \emph{Check-Effect Gap}. When an agent evaluates a precondition assertion on a mutable resource (e.g., verifying file permissions, checking inventory availability, or validating an account balance) at time $T_1$, and subsequently invokes a state-changing tool at time $T_2$, the state of the underlying resource may change between $T_1$ and $T_2$. This concurrency anomaly—analogous to classical Time-of-Check to Time-of-Use (TOCTOU) flaws in operating systems~\cite{bishop1996checking, wei2005toctou}—invalidates the precondition assertion upon which the agent's decision logic relied. In an agentic setting, this window can be exploited by adversarial external actors or triggered by asynchronous background processes, causing unauthorized mutations, resource corruption, or financial double-spending.

Classical program analysis and static taint-tracking techniques~\cite{arzt2014flowdroid, chen2002model} are fundamentally ill-equipped to detect check-effect binding risks in LLM agent workflows. Traditional taint analysis tracks data flows between sources and sinks but remains oblivious to semantic validity conditions: it cannot determine whether an intervening check operation actually binds its precondition to a downstream side effect. Furthermore, standard control-flow graphs fail to capture whether a tool invocation relies on authority-enforced concurrency mechanisms—such as atomic transactions, HTTP Entity Tags (ETags), conditional optimistic concurrency locks, or re-validation tokens.

To address these challenges, we introduce \textbf{AgentLint}, a contract-verified static and dynamic analysis framework designed specifically for detecting and repairing check-effect binding violations in LLM-agent tool workflows. The core insight behind AgentLint is that TOCTOU vulnerabilities in agent workflows are not merely temporal gaps, but structural \emph{binding failures} between precondition assertions and authority-enforced execution contracts. By elevating tool declarations into explicit semantic contracts that declare operation types, resource arguments, and binding protection mechanisms, AgentLint can statically reason about dependency graphs and track whether precondition validity is preserved across workflow execution paths.

We summarize the key contributions of this paper as follows:
\begin{itemize}
    \item \textbf{Formalization of Agent Check-Effect Bindings (\textsf{AWIR-B})}: We present a formal intermediate representation, the Agent Workflow IR with Bindings (\textsf{AWIR-B}), which models agent tool interactions as directed acyclic graphs equipped with explicit precondition assertions, resource targets, and authority-enforced protection mechanisms. We provide a relative soundness proof sketch for dependency propagation.
    \item \textbf{The AgentLint Framework}: We design and implement AgentLint, a static dependency propagation engine and finding generator that enforces topological reachability and contract-level binding checks. AgentLint produces prescriptive repair obligations to guide automated remediation.
    \item \textbf{The \textsf{BindBench} Benchmark Corpus}: We construct \textsf{BindBench}, a benchmark dataset comprising 24 distinct scenario families spanning cloud infrastructure, financial transactions, database locks, and filesystem identity. Each family includes vulnerable, repaired, and safe near-miss variants (72 workflows total), validated via dual independent human annotations with high inter-rater agreement ($\kappa \ge 0.70$).
    \item \textbf{Empirical Evaluation}: We evaluate AgentLint across six comprehensive experiments ($E1$–$E6$). AgentLint achieves perfect detection accuracy ($F_1 = 1.0$, $100\%$ precision and recall) on \textsf{BindBench}, completely eliminating the 48 false positives generated by reachability-only baselines while executing in under $1$\,ms per workflow. Dynamic validation confirms $100\%$ of flagged vulnerabilities under adversarial schedule mutation.
\end{itemize}

% TODO: [Future Work / Extension] Extend evaluation to open-source agent traces from wild MCP implementations (Experiment E8).
